\section{Implementation}
\label{sec:impl}

\tool is written in Python~3.14 in 10{,}034 lines across 38 files, grouped
into five packages: engine adapters, the corpus pipeline, the typed generator,
the oracle, and triage. It depends on four libraries: PyYAML for configuration,
requests for the GitHub API, sqlglot for reducing statements that parse, and
pytest. A suite of 514 tests (5{,}041 lines) passes in about 31 seconds. Two
problems we met while building it are worth reporting, because both produced
wrong results that looked like right ones.

\paragraph{Errors arrive in the middle of results.} \tool drives
\texttt{tursodb} as a subprocess rather than through a library binding. The
defects we look for end the process, and an in-process binding would take the
test harness down with it. The cost is a text protocol. Turso writes errors to
standard output, interleaved with result rows. Our first adapter looked for an
error only on the first line of the output. A query that returned two rows and
then failed therefore came back as three rows of data, the third being the error
message. The adapter now scans the whole output, and a regression test covers
this case.

\paragraph{A declared type is not a stored type.} The generator first assumed
that if \sql{CREATE TABLE t(c DATE)} parsed, the engine supported dates. SQLite
and Turso map any unknown type name to a storage class, so such a column really
holds integers. A generator built on that assumption reports type coverage it
does not have. \tool therefore probes the real binary once, records which types
it can tell apart and which functions exist, and refuses to emit anything else.

\paragraph{The seed is the state.} Every random choice comes from a seeded
generator, and a campaign writes its last seed on every progress report. A
campaign can stop and resume exactly, and any single case can be regenerated
from its seed alone. Section~\ref{sec:rq1} uses this to regenerate each defect.
